% from ega4-20-fr.tex
\oldpage[IV-4]{227}
$\vphi_{\beta\alpha}(S_\alpha)\subset S_\beta$ pour $\alpha\leq\beta$ ; on prend alors pour $S$ la limite inductive dans $A=\varinjlim A_\alpha$ du système inductif de parties $(S_\alpha)$.
\end{env}

\begin{env}[20.1.2]
\EGAReaderTarget{ega4.c04.statement.20-1-2}{IV.20.1.2}
Soit maintenant $\sh{F}$ un $\sh{O}_X$-module.
On pose alors
\[
\label{IV.20.1.2.1}
  \sh{F}[\sh{S}^{-1}]=\sh{F}\otimes_{\sh{O}_X}\sh{O}_X[\sh{S}^{-1}]
  \tag{20.1.2.1}
\]
et on dit que c'est le \emph{faisceau de modules de fractions de $\sh{F}$ à dénominateurs dans $\sh{S}$} ; il est immédiat qu'il est associé au préfaisceau de modules $U\mapsto\Gamma(U,\sh{F})[\Gamma(U,\sh{S})^{-1}]$, et que pour tout $x\in X$, on a un isomorphisme canonique
\[
\label{IV.20.1.2.2}
  (\sh{F}[\sh{S}^{-1}])_x\isoto\sh{F}_x[\sh{S}_x^{-1}].
  \tag{20.1.2.2}
\]
\end{env}

\begin{env}[20.1.3]
\EGAReaderTarget{ega4.c04.statement.20-1-3}{IV.20.1.3}
Nous allons nous intéresser ici au cas où $\sh{S}$ est le sous-faisceau $\sh{S}(\sh{O}_X)$ de $\sh{O}_X$ tel que pour tout ouvert $U$, $\Gamma(U,\sh{S})$ soit l'\emph{ensemble des éléments réguliers} de l'anneau $\Gamma(U,\sh{O}_X)$ ; il est immédiat qu'il s'agit d'un faisceau (et non seulement d'un préfaisceau), la régularité d'une section de $\sh{O}_X$ au-dessus de $U$ se vérifiant « fibre par fibre » (i.e. signifiant que le germe de la section en $x$ est régulier dans $\sh{O}_{X,x}$ pour tout $x\in U$), autrement dit $\sh{S}(\sh{O}_X)_x$ n'est autre que l'ensemble des éléments réguliers de $\sh{O}_{X,x}$.
Le faisceau d'anneaux correspondant
\[
  \sh{M}_X=\sh{O}_X[\sh{S}^{-1}]
\]
est appelé le \emph{faisceau des germes de fonctions méromorphes sur $X$}, et les sections de $\sh{M}_X$ au-dessus de $X$ sont appelées les \emph{fonctions méromorphes sur $X$} ; elles forment un anneau que l'on note $M(X)$.
Pour tout $\sh{O}_X$-module $\sh{F}$,
\[
  \sh{F}\otimes_{\sh{O}_X}\sh{M}_X=\sh{F}[\sh{S}^{-1}]
\]
est encore noté $\sh{M}_X(\sh{F})$ et appelé le \emph{faisceau des germes de sections méromorphes de $\sh{F}$} ; ses sections au-dessus de $X$ forment un $M(X)$-module noté $M(X,\sh{F})$, dont les éléments sont appelés \emph{sections méromorphes de $\sh{F}$ au-dessus de $X$}.
Ces définitions impliquent que pour tout ouvert $U$ de $X$, on a un isomorphisme canonique $\sh{M}_X(\sh{F})|U\isoto\sh{M}_U(\sh{F}|U)$, en particulier $\sh{M}_X|U\isoto\sh{M}_U$.
\end{env}

\begin{env}[20.1.3.1]
\label{IV.20.1.3.1}
Si $X$ est un \emph{préschéma réduit}, on notera que si un élément $s\in\Gamma(U,\sh{O}_X)$ est tel que $s_\xi\neq0$ pour tout point maximal $\xi$ de $U$, alors $s$ est \emph{régulier}.
En effet, si $st=0$ pour un $t\in\Gamma(U,\sh{O}_X)$, on a $s_\xi t_\xi=0$, donc $t_\xi=0$ puisque $\sh{O}_{X,\xi}$ est un corps, et dire que $t_\xi=0$ pour tout point maximal $\xi$ de $X$ signifie que $t=0$ : en effet, on est aussitôt ramené au cas où $U$ est affine, et un élément d'un anneau réduit qui appartient à tous les idéaux premiers minimaux est nul par définition.
La réciproque est vraie si l'ensemble des composantes irréductibles de $X$ est \emph{localement fini}.
On est en effet aussitôt ramené au cas où $X=\Spec(A)$ est affine ; si $\mathfrak p_i$ ($1\leq i\leq n$) sont les idéaux premiers minimaux de $A$ et si $s\in\mathfrak p_i$ pour un indice $i$, il existe $t\in A$ tel que $t\in\mathfrak p_j$ pour $j\neq i$ et $t\notin\mathfrak p_i$ (Bourbaki, \emph{Alg. comm.}, chap.~II, \textsection1, n\textsuperscript{o}1, prop.~1) ; on a donc $st\in\mathfrak p_i$ pour tout $i$, et par suite $st=0$ puisque $A$ est réduit ; $s$ est donc non régulier.
\end{env}

\begin{env}[20.1.4]
\EGAReaderTarget{ega4.c04.statement.20-1-4}{IV.20.1.4}
Pour tout ouvert $U$ de $X$, l'homomorphisme $t\mapsto t/1$ de $\Gamma(U,\sh{O}_X)$ dans $\Gamma(U,\sh{O}_X)[\Gamma(U,\sh{S})^{-1}]$ (qui n'est autre que l'\emph{anneau total des fractions} de
